\medskip

Le but de cet exercice est d'étudier la fonction $f$ définie sur l'intervalle $]0~;~+\infty[$ par : 

\[f(x) = x\ln \left(x^2\right) -\dfrac 1x.\]

\smallskip

\textbf{Partie A : lectures graphiques}

\medskip

On a tracé ci-dessous la courbe représentative $(\mathcal{C}_f)$ de la fonction $f$, ainsi que la droite $(T)$, tangente à la courbe $(\mathcal{C}_f)$ au point A de coordonnées $(1~;~-1)$.

Cette tangente passe également par le point B$(0~;~-4)$.

\begin{center}
\psset{xunit=1.5cm,yunit=0.6cm,arrowsize=2pt 3}
\begin{pspicture*}(-0.4,-7.2)(4,7.1)
\multido{\n=1+1}{4}{\psline[linewidth=0.2pt](\n,-7)(\n,7)}
\multido{\n=-7+1}{15}{\psline[linewidth=0.2pt](0,\n,)(4,\n)}
\psaxes[linewidth=1.25pt,labelFontSize=\scriptstyle,showorigin=false]{->}(0,0)(-0.4,-7)(4,7)
\psplot[plotpoints=2000,linewidth=1.25pt,linecolor=red]{0.1}{4}{x dup mul ln x mul 1 x div sub}
\psplotTangent{1}{3}{x dup mul ln x mul 1 x div sub}
\uput[l](3,6){\red $\mathcal{C}_f$}\uput[dr](3,5){$T$}
\uput[dr](1,-1){A}\uput[dr](0,-4){B}
\end{pspicture*}
\end{center}

\begin{enumerate}
\item Lire graphiquement $f'(1)$ et donner l'équation réduite de la tangente $(T)$.
\item Donner les intervalles sur lesquels la fonction $f$ semble convexe ou concave.

Que semble représenter le point A pour la courbe $(\mathcal{C}_f)$ ?
\end{enumerate}

\bigskip

\textbf{Partie B : étude analytique}

\medskip

\begin{enumerate}
\item Déterminer, en justifiant, la limite de $f$ en $+\infty$, puis sa limite en $0$.
\item On admet que la fonction $f$ est deux fois dérivable sur l'intervalle $]0~;~+\infty[$.
	\begin{enumerate}
		\item Déterminer $f'(x)$ pour $x$ appartenant à l'intervalle $]0~;~+\infty[$.
		\item Montrer que pour tout $x$ appartenant à l'intervalle $]0~;~+\infty[$,

\[f''(x) = \dfrac{2(x + 1)(x - 1)}{x^3}.\]
	\end{enumerate}
\item 
	\begin{enumerate}
		\item Étudier la convexité de la fonction $f$ sur l'intervalle $]0~;~+\infty[$.
		\item Étudier les variations de la fonction $f'$, puis le signe de $f'(x)$ pour $x$ appartenant à l'intervalle $]0~;~+\infty[$.
		
En déduire le sens de variation de la fonction $f$ sur l'intervalle $]0~;~+\infty[$.
	\end{enumerate}
\item
	\begin{enumerate}
		\item Montrer que l'équation $f(x) = 0$ admet une unique solution $\alpha$ sur l'intervalle $]0~;~+\infty[$.
		\item Donner la valeur arrondie au centième de $\alpha$ et montrer que $\alpha$ vérifie :
		
\[\alpha^2 = \text{exp} \left(\dfrac{1}{\alpha^2}\right).\]
	\end{enumerate}
\end{enumerate}

\bigskip
